% --- 题目 ---
\problem[P190-1 (25-3)]{设矩阵$\displaystyle \bm{A}=\begin{pmatrix} 1 & 2 \\ -2 & -a \end{pmatrix}$, $\displaystyle \bm{B}=\begin{pmatrix} 1 & 0 \\ 1 & a \end{pmatrix}$. 若$\displaystyle f(x,y)=|x\bm{A}+y\bm{B}|$是正定二次型,则$\displaystyle a$的取值范围是( \quad )
\begin{tasks}(2)
  \task $(0,2-\sqrt{3})$.
  \task $(2-\sqrt{3},2+\sqrt{3})$.
  \task $(2+\sqrt{3},4)$.
  \task $(0,4)$.
\end{tasks}}
\ansat{350；【十年真题】 - 考点三：正定二次型与正定矩阵 - 1}

% --- 题目 ---
\problem[P179-6 (96-1)]{设$\displaystyle \bm{A}=\bm{E}-\bm{\xi}\bm{\xi}^{\mathrm{T}}$, 其中$\displaystyle \bm{E}$是$\displaystyle n$阶单位矩阵, $\displaystyle \bm{\xi}$是$n$维非零列向量, $\displaystyle \bm{\xi}^{\mathrm{T}}$是$\displaystyle \bm{\xi}$的转置. 证明:
\begin{enumerate}
\item[(1)] $\displaystyle \bm{A}^2=\bm{A}$的充分条件是$\displaystyle \bm{\xi}^{\mathrm{T}}\bm{\xi}=1$
\item[(2)] 当$\displaystyle \bm{\xi}^{\mathrm{T}}\bm{\xi}=1$时, $\displaystyle \bm{A}$是不可逆矩阵.
\end{enumerate}}
\ansat{341；【真题精选】 - 考点：矩阵的特征值和特征向量 - 6}

% --- 题目 ---
\problem[P160-变式 (04-1)]{设有齐次线性方程组
\[ \begin{cases}
(1+a)x_1 + x_2 + \cdots + x_n = 0, \\
2x_1 + (2+a)x_2 + \cdots + 2x_n = 0, \\
\vdots \\
nx_1 + nx_2 + \cdots + (n+a)x_n = 0
\end{cases} (n \ge 2), \]
试问 $\displaystyle a$ 为何值时, 该方程组有非零解, 并求出其通解.}
\ansat{329；【方法探究】 - 考点一：线性方程组的解的情况及求解 - 变式}

% --- 题目 ---
\problem[P148-例2 (1)]{$\displaystyle n$ 阶行列式 $\displaystyle \begin{vmatrix} 1 & -1 & & & & \\ 2 & a & -1 & & & \\ 3 & & a & -1 & & \\ \vdots & & \ddots & \ddots & & \\ n-1 & & & & a & -1 \\ n & & & & & a \end{vmatrix} = \underline{\hspace{4em}}.$}
\ansat{148；【方法探究】 - 考点：具体行列式的计算 - 例2 (1)}

% --- 题目 ---
\problem[P180-4 (22-1)]{下述四个条件中, 3阶矩阵$\displaystyle \bm{A}$可对角化的一个充分但不必要条件是( \quad )
\begin{tasks}(2)
  \task $\displaystyle \bm{A}$有3个互不相等的特征值.
  \task $\displaystyle \bm{A}$有3个线性无关的特征向量.
  \task $\displaystyle \bm{A}$有3个两两线性无关的特征向量.
  \task $\displaystyle \bm{A}$的属于不同特征值的特征向量正交.
\end{tasks}}
\ansat{341；【十年真题】 - 考点：矩阵的相似和相似对角化 - 4}

% --- 题目 ---
\problem[P179-1 (08-1,2,3)]{设$\displaystyle \bm{A}$为$\displaystyle n$阶非零矩阵, $\displaystyle \bm{E}$为$\displaystyle n$阶单位矩阵. 若$\displaystyle \bm{A}^3=\bm{O}$,则( \quad )
\begin{tasks}(2)
  \task $\displaystyle \bm{E}-\bm{A}$不可逆, $\displaystyle \bm{E}+\bm{A}$不可逆.
  \task $\displaystyle \bm{E}-\bm{A}$不可逆, $\displaystyle \bm{E}+\bm{A}$可逆.
  \task $\displaystyle \bm{E}-\bm{A}$可逆, $\displaystyle \bm{E}+\bm{A}$可逆.
  \task $\displaystyle \bm{E}-\bm{A}$可逆, $\displaystyle \bm{E}+\bm{A}$不可逆.
\end{tasks}}
\ansat{341；【真题精选】 - 考点：矩阵的特征值和特征向量 - 1}

% --- 题目 ---
\problem[P170-4 (21-2)]{设 3 阶矩阵 $\displaystyle \bm{A}=(\bm{\alpha}_1, \bm{\alpha}_2, \bm{\alpha}_3), \ \bm{B}=(\bm{\beta}_1, \bm{\beta}_2, \bm{\beta}_3)$. 若向量组 $\displaystyle \bm{\alpha}_1, \bm{\alpha}_2, \bm{\alpha}_3$ 可以由向量组 $\displaystyle \bm{\beta}_1, \bm{\beta}_2, \bm{\beta}_3$ 线性表出, 则~(~\quad~)}
\begin{tasks}(2)
  \task $\displaystyle \bm{A}\bm{x}=\bm{0}$ 的解均为 $\displaystyle \bm{Bx}=\bm{0}$ 的解.
  \task $\displaystyle \bm{A}^{\mathrm{T}}\bm{x}=\bm{0}$ 的解均为 $\displaystyle \bm{B}^{\mathrm{T}}\bm{x}=\bm{0}$ 的解.
  \task $\displaystyle \bm{B}\bm{x}=\bm{0}$ 的解均为 $\displaystyle \bm{A}\bm{x}=\bm{0}$ 的解.
  \task $\displaystyle \bm{B}^{\mathrm{T}}\bm{x}=\bm{0}$ 的解均为 $\displaystyle \bm{A}^{\mathrm{T}}\bm{x}=\bm{0}$ 的解.
\end{tasks}
\ansat{336；【十年真题】 - 考点：线性方程组的解的结构 - 4}

% --- 题目 ---
\problem[P180-2 (24-2)]{设$\displaystyle \bm{A},\bm{B}$为2阶矩阵, 且$\displaystyle \bm{AB}=\bm{BA}$, 则“$\displaystyle \bm{A}$有两个不相等的特征值”是“$\displaystyle \bm{B}$可对角化”的( \quad )
\begin{tasks}(2)
  \task 充分必要条件.
  \task 充分不必要条件.
  \task 必要不充分条件.
  \task 既不充分也不必要条件.
\end{tasks}}
\ansat{341；【十年真题】 - 考点：矩阵的相似和相似对角化 - 2}

% --- 题目 ---
\problem[P159-1 (18-1,2,3)]{已知 $\displaystyle a$ 是常数, 且矩阵 $\displaystyle \bm{A} = \begin{pmatrix} 1 & 2 & a \\ 1 & 3 & 0 \\ 2 & 7 & -a \end{pmatrix}$ 可经初等列变换化为矩阵 $\displaystyle \bm{B} = \begin{pmatrix} 1 & a & 2 \\ 0 & 1 & 1 \\ -1 & 1 & 1 \end{pmatrix}.$}
\begin{enumerate}
\item[(1)] 求 $\displaystyle a$;
\item[(2)] 求满足 $\displaystyle \bm{AP} = \bm{B}$ 的可逆矩阵 $\displaystyle \bm{P}$.
\end{enumerate}
\begin{note}
  注意可逆这里的操作。
\end{note}
\ansat{328；【十年真题】 - 考点二：矩阵方程组的解的情况及求解 - 1}

% --- 题目 ---
\problem[P181-11 (18-2)]{设$\displaystyle \bm{A}$为3阶矩阵, $\displaystyle \bm{\alpha}_1,\bm{\alpha}_2,\bm{\alpha}_3$为线性无关的向量组. 若$\displaystyle \bm{A}\bm{\alpha}_1=2\bm{\alpha}_1+\bm{\alpha}_2+\bm{\alpha}_3$, $\displaystyle \bm{A}\bm{\alpha}_2=\bm{\alpha}_2+2\bm{\alpha}_3$, $\displaystyle \bm{A}\bm{\alpha}_3=-\bm{\alpha}_2+\bm{\alpha}_3$,则$\bm{A}$的实特征值为\underline{\hspace{4em}}.}
\ansat{342；【十年真题】 - 考点：矩阵的相似和相似对角化 - 11}

% --- 题目 ---
\problem[P173-10 (05-1,2)]{已知3阶矩阵$\displaystyle \bm{A}$的第一行是$\displaystyle (a,b,c)$, $\displaystyle a,b,c$不全为零,矩阵 $\displaystyle \bm{B}=\begin{pmatrix} 1 & 2 & 3 \\ 2 & 4 & 6 \\ 3 & 6 & k \end{pmatrix}$ ($\displaystyle k$为常数), 且$\displaystyle \bm{AB}=\bm{O}$, 求线性方程组$\displaystyle \bm{A}\bm{x}=\bm{0}$的通解.}
\ansat{338；【真题精选】 - 考点：线性方程组的解的结构 - 10}

% --- 题目 ---
\problem[P163-2 (23-1)]{已知 $\displaystyle n$ 阶矩阵 $\displaystyle \bm{A}, \bm{B}, \bm{C}$ 满足 $\displaystyle \bm{ABC} = \bm{O}, \bm{E}$ 为 $\displaystyle n$ 阶单位矩阵. 记矩阵
$$ 
\begin{pmatrix} \bm{O} & \bm{A} \\ \bm{BC} & \bm{E} \end{pmatrix}, \; \begin{pmatrix} \bm{AB} & \bm{C} \\ \bm{O} & \bm{E} \end{pmatrix}, \; \begin{pmatrix} \bm{E} & \bm{AB} \\ \bm{AB} & \bm{O} \end{pmatrix}
$$ 
的秩分别为 $\displaystyle r_1, r_2, r_3$, 则~(~\quad~)}
\begin{tasks}(2)
  \task $\displaystyle r_1 \le r_2 \le r_3.$
  \task $\displaystyle r_1 \le r_3 \le r_2.$
  \task $\displaystyle r_3 \le r_1 \le r_2.$
  \task $\displaystyle r_2 \le r_1 \le r_3.$
\end{tasks}
\ansat{333；【十年真题】 - 考点：向量组的线性相关性、线性表示及秩 - 2}

% --- 题目 ---
\problem[P150-3 (24-1)]{设实矩阵 $\displaystyle \bm{A} = \begin{pmatrix} a+1 & a \\ a & a \end{pmatrix}$. 若对任意实向量 $\displaystyle \boldsymbol{\alpha} = \begin{pmatrix} x_1 \\ x_2 \end{pmatrix}, \ \boldsymbol{\beta} = \begin{pmatrix} y_1 \\ y_2 \end{pmatrix}$, 
$$\displaystyle (\boldsymbol{\alpha}^{\mathrm{T}} \bm{A} \boldsymbol{\beta})^2 \le \boldsymbol{\alpha}^{\mathrm{T}} \bm{A} \boldsymbol{\alpha} \cdot \boldsymbol{\beta}^{\mathrm{T}} \bm{A} \boldsymbol{\beta}$$ 
都成立, 则 $\displaystyle a$ 的取值范围为 \underline{\hspace{4em}}.}
\ansat{325；【十年真题】 - 考点一：矩阵的运算 - 3}

% --- 题目 ---
\problem[P146-4 (20-1,2,3)]{行列式 $\displaystyle \begin{vmatrix} a & 0 & -1 & 1 \\ 0 & a & 1 & -1 \\ -1 & 1 & a & 0 \\ 1 & -1 & 0 & a \end{vmatrix} = \underline{\hspace{4em}}.$}
\ansat{323；【十年真题】 - 考点：具体行列式的计算 - 4}

% --- 题目 ---
\problem[P180-5 (22-2,3)]{设$\displaystyle \bm{A}$为3阶矩阵, $\displaystyle \bm{\Lambda}=\begin{pmatrix} 1 & 0 & 0 \\ 0 & -1 & 0 \\ 0 & 0 & 0 \end{pmatrix}$, 则$\displaystyle \bm{A}$的特征值为 $\displaystyle 1,-1,0$ 的充分必要条件是( \quad )
\begin{tasks}(2)
  \task 存在可逆矩阵$\displaystyle \bm{P},\bm{Q}$, 使得$\displaystyle \bm{A}=\bm{P}\bm{\Lambda}\bm{Q}$.
  \task 存在可逆矩阵$\displaystyle \bm{P}$, 使得$\displaystyle  \bm{A}=\bm{P}\bm{\Lambda}\bm{P}^{-1}$.
  \task 存在正交矩阵$\displaystyle \bm{Q}$, 使得$\displaystyle \bm{A}=\bm{Q}\bm{\Lambda}\bm{Q}^{-1}$.
  \task 存在可逆矩阵$\displaystyle \bm{P}$, 使得$\displaystyle \bm{A}=\bm{P}\bm{\Lambda}\bm{P}^{\mathrm{T}}$.
\end{tasks}}
\ansat{341；【十年真题】 - 考点：矩阵的相似和相似对角化 - 5}

% --- 题目 ---
\problem[P186-10 (02-1)]{设$\displaystyle \bm{A}, \bm{B}$为同阶方阵.
\begin{enumerate}
\item[(1)] 如果$\displaystyle \bm{A}, \bm{B}$相似, 试证$\displaystyle \bm{A}, \bm{B}$的特征多项式相等;
\end{enumerate}}
\ansat{345；【真题精选】 - 考点：矩阵的相似和相似对角化 - 10}

% --- 题目 ---
\problem[P156-10 (95-1)]{设 $\displaystyle \bm{A}$ 是 $\displaystyle n$ 阶矩阵, 满足 $\displaystyle \bm{A}\bm{A}^{\mathrm{T}}=\bm{E}$ ($\displaystyle \bm{E}$ 是 $\displaystyle n$ 阶单位矩阵, $\displaystyle \bm{A}^{\mathrm{T}}$ 是 $\displaystyle \bm{A}$ 的转置矩阵), $\displaystyle |\bm{A}| < 0$, 则 $\displaystyle |\bm{A} + \bm{E}| = \underline{\hspace{4em}}. $}
\ansat{327；【真题精选】 - 考点一：矩阵的运算 - 10}

% --- 题目 ---
\problem[P149-4 (96-5)]{5 阶行列式 $\displaystyle \begin{vmatrix} 1-a & a & 0 & 0 & 0 \\ -1 & 1-a & a & 0 & 0 \\ 0 & -1 & 1-a & a & 0 \\ 0 & 0 & -1 & 1-a & a \\ 0 & 0 & 0 & -1 & 1-a \end{vmatrix} = \underline{\hspace{4em}}.$}
\ansat{324；【真题精选】 - 考点：具体行列式的计算 - 4}
